\documentclass{article}
\usepackage{amsmath, amssymb, amsthm}
\usepackage[margin=1in]{geometry}

\newcommand{\E}{\mathbb{E}}
\newcommand{\Prob}{\mathbb{P}}

\newtheorem{definition}{Definition}
\newtheorem{proposition}{Proposition}

\title{Threshold Strategy: Unifying Early and Late Arrival}
\author{}
\date{}

\begin{document}
\maketitle

\section{Motivation}

The early-arrival strategy maximizes lap count by minimizing wait time. The late-arrival strategy maximizes lifespan utilization by positioning carefully. Neither is universally optimal.

A \textbf{threshold strategy} adapts based on the current lifespan $l$:
\begin{itemize}
    \item Short $l$: Re-roll quickly (early behavior)
    \item Long $l$: Position strategically (late behavior)
\end{itemize}

This document explores what the threshold function $L^*(\Omega)$ should look like.

\section{Setup}

At state $(\Omega, l)$, the particle can:
\begin{itemize}
    \item \textbf{Go early}: Wait $w_E$ (minimal), arrive at $\Omega_E$, cost $w_E + \tau$
    \item \textbf{Go late}: Wait $w_L$ (maximal), arrive at $\Omega_L = T_1$, cost $w_L + \tau$
\end{itemize}

Feasibility constraints:
\begin{align}
    \text{Early feasible} &\iff w_E + \tau \leq l \\
    \text{Late feasible} &\iff w_L + \tau \leq l
\end{align}

Since $w_E \leq w_L$, if late is feasible then early is also feasible.

\section{Deriving the Threshold}

\subsection{Attempt 1: Feasibility-Based Threshold}

The simplest threshold: go late when you can, early otherwise.
\[
L^*_{\text{feasibility}}(\Omega) = w_L(\Omega) + \tau
\]

This is what the early strategy already does as a fallback. It doesn't really unify the strategies.

\subsection{Attempt 2: Expected Value Comparison}

Compare single-lap expected values:
\begin{align}
    V_E &= 1 + U(\Omega_E) \quad \text{(early)} \\
    V_L &= 1 + U(\Omega_L) \quad \text{(late)}
\end{align}

Go late if $V_L > V_E$, i.e., if $U(\Omega_L) > U(\Omega_E)$.

But this ignores the \emph{cost} difference. Late uses more lifespan ($w_L + \tau$ vs $w_E + \tau$). We should account for opportunity cost.

\subsection{Attempt 3: Rate-Based Comparison}

Consider the ``value rate'' — expected future laps per unit time spent:
\begin{align}
    R_E &= \frac{1 + U(\Omega_E)}{w_E + \tau} \\
    R_L &= \frac{1 + U(\Omega_L)}{w_L + \tau}
\end{align}

Go late if $R_L > R_E$. This accounts for the time cost of waiting.

But this still doesn't depend on $l$. We want the threshold to use information about our current lifespan.

\subsection{Attempt 4: Lifespan-Aware Threshold}

The key insight: with a long lifespan, you can afford to position carefully. With a short lifespan, you need to act fast.

Consider: if $l$ is large, you might be able to do \emph{multiple} early laps before needing to worry. The question is whether multiple early laps beats one late lap.

Let $k = \lfloor (l - \tau) / \tau \rfloor$ be the maximum number of immediate laps possible (ignoring gate constraints). The value of ``spamming'' $k$ early laps is roughly:
\[
V_{\text{spam}} \approx k + U(\Omega + k\tau)
\]

Compare to one late lap:
\[
V_{\text{late}} = 1 + U(T_1)
\]

The threshold should be where these are equal. Roughly:
\[
L^* \approx \tau \cdot \left( U(T_1) - U(\Omega_E) + 1 \right)
\]

This suggests: if the positional advantage $U(T_1) - U(\Omega_E)$ is large, set a high threshold (favor late). If it's small, set a low threshold (favor early).

\subsection{Attempt 5: Simple Quantile-Based Threshold}

A pragmatic approach: set the threshold at some quantile of the lifespan distribution.

Let $L^* = c / \lambda$ for some constant $c > 0$. Then:
\begin{itemize}
    \item $c = 0$: Always late (when feasible)
    \item $c = \infty$: Always early
    \item $c = 1$: Threshold at the mean lifespan
    \item $c = \ln 2 \approx 0.69$: Threshold at the median
\end{itemize}

The optimal $c$ likely depends on the velocity regime:
\begin{itemize}
    \item Fast ($\tau \ll T_1$): High $c$ (favor early, spam laps)
    \item Slow ($\tau > T_1$): Low $c$ (favor late, position matters)
\end{itemize}

\section{Proposed Threshold Function}

Combining the above insights, a reasonable threshold is:
\[
L^*(\Omega) = \alpha \cdot \frac{1}{\lambda} + \beta \cdot (w_L(\Omega) - w_E(\Omega))
\]

where:
\begin{itemize}
    \item $\alpha \geq 0$ controls the baseline threshold (as fraction of mean lifespan)
    \item $\beta \geq 0$ scales with the wait-time difference (cost of going late)
\end{itemize}

Special cases:
\begin{itemize}
    \item $\alpha = 0, \beta = 0$: Always late (if feasible)
    \item $\alpha = \infty$: Always early
    \item $\alpha = 1, \beta = 0$: Go early if $l < $ mean, late if $l \geq$ mean
\end{itemize}

\section{Simplified Threshold for Implementation}

For initial implementation, use a single-parameter threshold:
\[
L^*(\Omega) = c \cdot \frac{1}{\lambda}
\]

where $c$ is a tunable constant. The strategy becomes:
\[
w^* = \begin{cases}
    w_E(\Omega, l) & \text{if } l < L^*(\Omega) \\
    w_L(\Omega, l) & \text{if } l \geq L^*(\Omega) \text{ and late is feasible} \\
    w_E(\Omega, l) & \text{otherwise (fallback)}
\end{cases}
\]

We can then sweep over $c \in [0, 5]$ to find the optimal threshold empirically.

\section{Empirical Results}

We implemented the threshold strategy and swept over $c$ values in two regimes.

\subsection{Fast Regime: $\tau = 1 < T_1 = 2$}

Parameters: $v = 2$, $\tau = 1$, $T_1 = 2$, $T_2 = 1$, mean lifespan $= 3$ ($\lambda = 1/3$), $\Omega = 0.5$.

\begin{center}
\begin{tabular}{|c|c|}
\hline
\textbf{Strategy} & \textbf{Expected Laps} \\
\hline
Late & 1.337 \\
Threshold $c = 0.5$ & 1.388 \\
Threshold $c = 1.0$ & 1.843 \\
Threshold $c = 2.0$ & 1.893 \\
Threshold $c = 3.0$ & 1.900 \\
Early & 1.903 \\
\hline
\end{tabular}
\end{center}

\textbf{Finding}: In the fast regime, pure early is optimal. As $c \to \infty$, the threshold strategy converges to early. There is no benefit to the hybrid approach.

\subsection{Slow Regime: $\tau = 2.5 > T_1 = 2$}

Parameters: $v = 0.8$, $\tau = 2.5$, $T_1 = 2$, $T_2 = 1$, mean lifespan $= 3$ ($\lambda = 1/3$), $\Omega = 0.5$.

\begin{center}
\begin{tabular}{|c|c|}
\hline
\textbf{Strategy} & \textbf{Expected Laps} \\
\hline
Early & 0.605 \\
Threshold $c = 2.0$ & 0.628 \\
Late & 0.631 \\
Threshold $c = 0.5$ & 0.631 \\
Threshold $c = 1.5$ & 0.635 \\
\textbf{Threshold $c = 1.0$} & \textbf{0.638} \\
\hline
\end{tabular}
\end{center}

\textbf{Finding}: In the slow regime, the threshold strategy with $c \approx 1$ (threshold at mean lifespan) \emph{beats both pure strategies}. This confirms the hypothesis that a hybrid approach can outperform either extreme.

\section{Interpretation}

The results align with our intuition:

\textbf{Fast regime} ($\tau < T_1$): Particles can complete multiple laps per ON window. The ``spam laps'' effect dominates—rapid re-rolling of lifespans is always beneficial. Positioning provides negligible advantage because survival per lap is nearly certain.

\textbf{Slow regime} ($\tau > T_1$): Each lap necessarily crosses the OFF period. Here, the threshold strategy provides genuine value:
\begin{itemize}
    \item \textbf{Short lifespan} ($l < $ mean): The particle is unlikely to survive long anyway. Use the early strategy to attempt a quick lap and re-roll.
    \item \textbf{Long lifespan} ($l \geq $ mean): The particle has ``won the lottery.'' Use this valuable lifespan strategically by positioning at $T_1$ for optimal future phase.
\end{itemize}

The optimal threshold at $c = 1$ (the mean) makes intuitive sense: lifespans below average are ``unlucky'' and should be re-rolled; lifespans above average are ``lucky'' and should be exploited.

\section{Conclusions}

\begin{enumerate}
    \item The threshold strategy unifies early and late as special cases ($c = \infty$ and $c = 0$ respectively).
    
    \item In the fast regime, pure early is optimal. The threshold strategy offers no improvement.
    
    \item In the slow regime, the threshold strategy with $c \approx 1$ outperforms both pure strategies, achieving a $\sim$1\% improvement over late and $\sim$5\% improvement over early.
    
    \item The optimal threshold is remarkably simple: $L^* = \text{mean lifespan} = 1/\lambda$.
\end{enumerate}

\section{Open Questions}

\begin{enumerate}
    \item Does the optimal $c$ depend on $\Omega$? Our tests used fixed $\Omega = 0.5$.
    
    \item Is there a velocity ``crossover'' regime where the optimal $c$ transitions smoothly from high to low?
    
    \item Can we derive $c^* = 1$ analytically from the MDP structure?
    
    \item For maximizing \emph{survival time} rather than lap count, is late (or a different threshold) optimal?
\end{enumerate}

\end{document}
